\chapter{Backend}

Il backend di Ruggine è sviluppato in Rust per garantire performance, sicurezza e concorrenza. L'adozione di Axum, SQLx e Tokio consente di gestire richieste HTTP/WebSocket in modo asincrono e scalabile. La struttura a livelli separa chiaramente responsabilità e facilita la manutenzione.

\section{Struttura interna}
Il backend segue il pattern controller-service-repository. I controller gestiscono le richieste in ingresso (HTTP\/WS), validano i dati e orchestrano la logica applicativa; ad esempio, authentication\_controller.rs gestisce login e registrazione. I service implementano la business logic, come validazione, hashing password, generazione token, gestione gruppi; auth\_service.rs contiene le funzioni register\_user e login\_user. I repository si occupano dell'accesso ai dati, eseguendo query sicure e ottimizzate; user\_repository.rs offre funzioni CRUD sugli utenti.

\subsection{Esempio: flusso di registrazione}
Il flusso di registrazione prevede che il controller riceva la richiesta e chiami la funzione register\_user del service. Il service valida i dati, esegue l'hashing della password e chiama il repository, che salva l'utente nel database. Infine, il controller restituisce una risposta JSON.

\begin{lstlisting}[language=Rust, caption={Estratto da authentication\_controller.rs}]
pub async fn register_handler(State(state): State<Arc<AppState>>, Json(payload): Json<AuthPayload>) -> Result<Json<AuthResponse>, AppError> {
    auth_service::register_user(&state.db, payload).await?;
    Ok(Json(AuthResponse { message: "Utente registrato con successo".to_string() }))
}
\end{lstlisting}

\section{Funzioni principali e motivazioni}
	\textbf{register\_user}: valida input, controlla unicità username, esegue hashing sicuro (Argon2), salva utente e lo aggiunge al gruppo globale.\newline
	\textbf{login\_user}: valida input, verifica credenziali, genera JWT token per autenticazione sicura.\newline
	\textbf{create\_user/find\_by\_username} (repository): funzioni CRUD ottimizzate e sicure contro SQL injection.

\section{Funzioni principali e motivazioni}
\subsection{Registrazione utente: register\_user}
Questa funzione valida i dati ricevuti, controlla che l'username sia unico, esegue l'hashing sicuro della password con Argon2 e salva l'utente nel database. Infine aggiunge l'utente al gruppo globale.\newline
	extbf{Motivazione:} garantire sicurezza e integrit\`a dei dati gi\`a in fase di registrazione.
\begin{lstlisting}[language=Rust, caption={Estratto da auth\_service.rs - register\_user}]
pub async fn register_user(db: &SqlitePool, payload: AuthPayload) -> Result<(), AppError> {
    // Validazione
    if payload.username.trim().is_empty() { ... }
    // Hashing password
    let password_hash = hash_password(&payload.password)?;
    // Controllo unicita'
    if user_repository::find_by_username(db, &payload.username).await?.is_some() { ... }
    // Salvataggio
    user_repository::create_user(db, &payload.username, &password_hash).await?;
    // Aggiunta al gruppo globale
    group_repository::add_member_or_ignore(db, 1, new_user.id).await?;
    Ok(())
}
\end{lstlisting}

\subsection{Login utente: login\_user}
Valida le credenziali, verifica la password tramite Argon2 e genera un JWT token per l'autenticazione sicura lato client.\newline
	extbf{Motivazione:} autenticazione robusta e stateless.
\begin{lstlisting}[language=Rust, caption={Estratto da auth\_service.rs - login\_user}]
pub async fn login_user(db: &SqlitePool, payload: AuthPayload) -> Result<String, AppError> {
    // Validazione input
    // Ricerca utente
    // Verifica password
    // Generazione JWT
    Ok(token)
}
\end{lstlisting}

\subsection{Invio messaggio: save\_message}
Gestisce la ricezione di un messaggio da WebSocket, lo valida e lo salva nella tabella messages, poi lo inoltra agli altri membri del gruppo.\newline
	extbf{Motivazione:} garantire persistenza e distribuzione in tempo reale.
\begin{lstlisting}[language=Rust, caption={Flusso invio messaggio in chat}]
pub async fn save_message(db: &SqlitePool, group_id: i64, sender_id: i64, content: &str) -> Result<(), AppError> {
    // Validazione
    // Salvataggio su DB
    // Notifica agli altri membri
    Ok(())
}
\end{lstlisting}

\subsection{Gestione inviti: invitation\_repository}
Permette di creare, accettare o rifiutare inviti a gruppi. Ogni operazione aggiorna la tabella invitations e notifica gli utenti coinvolti.\newline
	extbf{Motivazione:} gestione collaborativa e sicura dei gruppi.
\begin{lstlisting}[language=Rust, caption={Esempio gestione invito}]
pub async fn create_invitation(db: &SqlitePool, from: i64, to: i64, group: i64) -> Result<(), AppError> {
    // Inserimento invito
    Ok(())
}
\end{lstlisting}

\section{Gestione delle richieste e flussi}
\begin{itemize}
    \item \textbf{HTTP}: usato per autenticazione, invio inviti, gestione utenti.
    \item \textbf{WebSocket}: usato per chat in tempo reale, notifiche, aggiornamenti gruppi/inviti.
\end{itemize}
Il controller WebSocket gestisce la connessione, autentica tramite token e smista i messaggi ai service appropriati.

\subsection{Esempio: invio messaggio in chat}
\begin{lstlisting}[language=Rust, caption={Handler WebSocket per messaggi}]
// Riceve messaggio, valida, salva e inoltra
pub async fn handle_group_message(...) -> Result<(), AppError> {
    chat_service::save_message(...).await?;
    // Inoltra ai membri del gruppo
}
\end{lstlisting}

\section{Sicurezza e validazione}
\begin{itemize}
    \item Validazione input centralizzata nei service.
    \item Hashing password con Argon2.
    \item Generazione e verifica JWT per autenticazione.
    \item Gestione errori tramite modulo dedicato (\texttt{error.rs}).
\end{itemize}

\section{Scelte progettuali}
Rust garantisce sicurezza a livello di memoria e performance elevate. L'architettura a livelli isola le responsabilità e semplifica l'estensione. L'uso di WebSocket permette una chat reattiva e scalabile. La modularità facilita test, manutenzione e aggiunta di nuove feature.